% ef-light-cones.tex — Light cones in Schwarzschild and ingoing EF coordinates
% Inputable TikZ fragment for fig:ef-light-cones
% Requires: tikz with calc, arrows.meta (loaded by ehbook.cls)
%
% Two-panel comparison showing how light cones tilt at the event horizon.
% Panel (a): Schwarzschild (r, t) — symmetric cones pinching at r=2M.
% Panel (b): Ingoing EF (r, v)  — ingoing edge horizontal, outgoing tilts inward.
%
% Physics (M=1): null directions from ds²=0.
%   (a) dt/dr = ±1/(1−2/r).  Symmetric about time axis.
%   (b) Ingoing: dv=0.  Outgoing: dv/dr = 2/(1−2/r).
%
% Scales: s_r = 0.7 cm/M, s_t = s_v = 0.35 cm/M.  Arm length L = 0.35 cm.
%
\begin{tikzpicture}[
    >={Stealth[length=5pt]},
    font=\footnotesize,
  ]

  % ── Panel (a): Schwarzschild coordinates ───────────────────────────

  \begin{scope}[shift={(0,0)}]

    % Axis lengths
    \def\rmin{0.5}    % left edge in cm (r ≈ 0.7M)
    \def\rmax{3.5}    % right edge in cm (r ≈ 5M)
    \def\tmax{2.6}    % top of t axis

    % Horizon position: r=2M → x = 2×0.7 = 1.4 cm
    \def\rh{1.4}

    % Gray shaded strip for r < 2M (coordinates undefined)
    \fill[EHMarginGray, opacity=0.06]
      (0, -0.3) rectangle (\rh, \tmax);

    % Horizon line
    \draw[EHHorizonCrimson, dashed, thick]
      (\rh, -0.3) -- (\rh, \tmax)
      node[above, font=\scriptsize, text=EHHorizonCrimson] {$r = 2M$};

    % Axes
    \draw[->, EHMarginGray, thin] (-0.15, 0) -- (\rmax, 0)
      node[right, font=\scriptsize] {$r$};
    \draw[->, EHMarginGray, thin] (0, -0.15) -- (0, \tmax)
      node[above, font=\scriptsize] {$t$};

    % ── Light cones ──────────────────────────────────────────────────
    % Each cone: filled triangle with base at (r, 0), edges ±(Δx, Δy).
    % Pre-computed from plan: dt/dr = ±1/f, f = 1−2/r, arm length 0.35 cm.

    % r = 4.2M → x = 2.94
    \fill[EHAccretionGold, opacity=0.15]
      (2.94, 0) -- ++(0.253, 0.242) -- ++(-0.506, 0) -- cycle;
    \draw[EHAccretionGold, thick]
      (2.94, 0) -- ++(0.253, 0.242)  (2.94, 0) -- ++(-0.253, 0.242);
    \fill[EHAccretionGold] (2.94, 0) circle (1.2pt);

    % r = 3.2M → x = 2.24
    \fill[EHAccretionGold, opacity=0.15]
      (2.24, 0) -- ++(0.210, 0.280) -- ++(-0.420, 0) -- cycle;
    \draw[EHAccretionGold, thick]
      (2.24, 0) -- ++(0.210, 0.280)  (2.24, 0) -- ++(-0.210, 0.280);
    \fill[EHAccretionGold] (2.24, 0) circle (1.2pt);

    % r = 2.6M → x = 1.82
    \fill[EHAccretionGold, opacity=0.15]
      (1.82, 0) -- ++(0.147, 0.318) -- ++(-0.294, 0) -- cycle;
    \draw[EHAccretionGold, thick]
      (1.82, 0) -- ++(0.147, 0.318)  (1.82, 0) -- ++(-0.147, 0.318);
    \fill[EHAccretionGold] (1.82, 0) circle (1.2pt);

    % r = 2.2M → x = 1.54
    \fill[EHAccretionGold, opacity=0.15]
      (1.54, 0) -- ++(0.063, 0.344) -- ++(-0.126, 0) -- cycle;
    \draw[EHAccretionGold, thick]
      (1.54, 0) -- ++(0.063, 0.344)  (1.54, 0) -- ++(-0.063, 0.344);
    \fill[EHAccretionGold] (1.54, 0) circle (1.2pt);

    % Panel label
    \node[font=\scriptsize, anchor=north] at (1.75, -0.55) {(a) Schwarzschild};

  \end{scope}

  % ── Panel (b): Ingoing Eddington–Finkelstein coordinates ──────────

  \begin{scope}[shift={(5.0, 0)}]

    % Axis lengths
    \def\rmin{0.2}
    \def\rmax{3.5}
    \def\vmax{2.6}

    % Horizon position: r=2M → x = 1.4 cm
    \def\rh{1.4}

    % Horizon line
    \draw[EHHorizonCrimson, dashed, thick]
      (\rh, -0.3) -- (\rh, \vmax)
      node[above, font=\scriptsize, text=EHHorizonCrimson] {$r = 2M$};

    % Axes
    \draw[->, EHMarginGray, thin] (-0.15, 0) -- (\rmax, 0)
      node[right, font=\scriptsize] {$r$};
    \draw[->, EHMarginGray, thin] (0, -0.15) -- (0, \vmax)
      node[above, font=\scriptsize] {$v$};

    % ── Light cones ──────────────────────────────────────────────────
    % Ingoing edge: dv=0, always (−0.35, 0) — horizontal pointing left.
    % Outgoing edge: dv/dr = 2/(1−2/r), pre-computed per plan.

    % r = 4.2M → x = 2.94
    \fill[EHAccretionGold, opacity=0.15]
      (2.94, 0) -- ++(-0.350, 0) -- ++(0.513, 0.310) -- cycle;
    \draw[EHAccretionGold, thick]
      (2.94, 0) -- ++(-0.350, 0)  (2.94, 0) -- ++(0.163, 0.310);
    \fill[EHAccretionGold] (2.94, 0) circle (1.2pt);

    % r = 3.2M → x = 2.24
    \fill[EHAccretionGold, opacity=0.15]
      (2.24, 0) -- ++(-0.350, 0) -- ++(0.473, 0.328) -- cycle;
    \draw[EHAccretionGold, thick]
      (2.24, 0) -- ++(-0.350, 0)  (2.24, 0) -- ++(0.123, 0.328);
    \fill[EHAccretionGold] (2.24, 0) circle (1.2pt);

    % r = 2.6M → x = 1.82
    \fill[EHAccretionGold, opacity=0.15]
      (1.82, 0) -- ++(-0.350, 0) -- ++(0.429, 0.341) -- cycle;
    \draw[EHAccretionGold, thick]
      (1.82, 0) -- ++(-0.350, 0)  (1.82, 0) -- ++(0.079, 0.341);
    \fill[EHAccretionGold] (1.82, 0) circle (1.2pt);

    % r = 2.0M → x = 1.40 (at the horizon — outgoing is vertical)
    \fill[EHAccretionGold, opacity=0.15]
      (1.40, 0) -- ++(-0.350, 0) -- ++(0.350, 0.350) -- cycle;
    \draw[EHAccretionGold, thick]
      (1.40, 0) -- ++(-0.350, 0)  (1.40, 0) -- ++(0, 0.350);
    \fill[EHAccretionGold] (1.40, 0) circle (1.2pt);

    % r = 1.5M → x = 1.05 (inside — outgoing tilts inward)
    \fill[EHAccretionGold, opacity=0.15]
      (1.05, 0) -- ++(-0.350, 0) -- ++(0.239, 0.332) -- cycle;
    \draw[EHAccretionGold, thick]
      (1.05, 0) -- ++(-0.350, 0)  (1.05, 0) -- ++(-0.111, 0.332);
    \fill[EHAccretionGold] (1.05, 0) circle (1.2pt);

    % r = 1.0M → x = 0.70 (deep inside — both edges point left)
    \fill[EHAccretionGold, opacity=0.15]
      (0.70, 0) -- ++(-0.350, 0) -- ++(0.103, 0.247) -- cycle;
    \draw[EHAccretionGold, thick]
      (0.70, 0) -- ++(-0.350, 0)  (0.70, 0) -- ++(-0.247, 0.247);
    \fill[EHAccretionGold] (0.70, 0) circle (1.2pt);

    % Annotation: arrow showing both edges point to smaller r
    \draw[EHMarginGray, thin, ->]
      (0.25, 0.85) -- ++(-0.20, 0)
      node[left, font=\scriptsize, text=EHMarginGray, align=right]
      {both edges\\[-1pt]point inward};

    % Panel label
    \node[font=\scriptsize, anchor=north] at (1.75, -0.55) {(b) Eddington--Finkelstein};

  \end{scope}

\end{tikzpicture}
